Regularized multi-task learning (MTL) is usually judged against independent per-task models. On the School benchmark (139 schools, 15,362 students), the opposite extreme, a single model pooled across all schools, has mostly been reported under a 75\% training protocol, with conflicting conclusions: it was the best method of \citet{evgeniou2005kernel} and clearly behind multi-task feature learning in \citet{argyriou2008convex}. We revisit this comparison across training fractions, with a focus on the low-data regime. A short NumPy implementation first reproduces published results, matching the trace-norm explained variance of \citet{argyriou2008convex} (26.5\% against 26.7\%). Across training fractions from 5\% to 75\% of each school, trace-norm MTL beats independent ridge on every split, by 19.2\% in normalized mean squared error at 5\% and 4.4\% at 75\%, yet the pooled model beats it on every split up to 20\%. Two facts explain the reversal. School-level covariates are constant within each school, so their effects are identifiable only across schools, and the trace-norm and $\ell_{2,1}$ penalties shrink every school towards zero: at 5\% the average trace-norm school model has 57\% of the norm of the pooled model. Mean-regularized MTL, which learns a shared mean and penalizes only deviations from it, stays within 0.4\% of pooled ridge at 5\%, beats it on every split from 10\% onwards by up to 6.3\%, and has lower error than trace-norm MTL at every training fraction; constraining the deviations to be low rank adds nothing on this benchmark. We also show that a fixed iteration budget fails silently on unscaled features, and that two common definitions of explained variance differ by 11 to 13 points for identical predictions.
